\section{坐标系为何是 Part-Whole Relation 的核心}

如果一个 nostril 支持“这里有一个 nose”，它不仅要预测 identity，还要预测 nose 的 position、orientation 与 scale。换句话说，part-whole relation 不是一个离散标签，而是 intrinsic coordinate frame 之间的变换。Hinton 用一组心理实验先让听众感受到：我们对同一组线条的理解，会随内部 reference frame 改变。

\lecturefigure{slide-08-psychological-coordinate-frames.jpg}{接下来几页用心理演示说明视觉 parse 与坐标系的真实性。}{Stanford Online 官方视频 00:09:16--00:13:31。}

\subsection{Cube demonstration：难点不是旋转，而是换 frame}

课堂要求听众想象一个 wire-frame cube，把原来的 body diagonal 旋转到竖直方向，再指出其他角的位置。许多人并不是不会三维几何，而是习惯把 cube 放在“桌面水平、竖直边朝上”的 intrinsic frame 中；一旦强制使用陌生 frame，同一对象的 internal representation 会剧烈改变。

\lecturefigure{slide-09-cube-demonstration.jpg}{Cube mental-rotation task：把熟悉的 body diagonal 改成竖直轴。}{Stanford Online 官方视频 00:13:32--00:16:07。}

\begin{knowledgebox}{读图：任务真正改变的是 intrinsic frame}
先找原来接触桌面的 front-bottom-right corner，再找与它相对的 top-back-left corner；课堂要求把这条 body diagonal 旋成竖直轴。其他六个 corner 的几何关系没有消失，但熟悉的“底面水平、竖边朝上”坐标系被撤走。若听众难以定位角点，证据支持的是 reference-frame dependence，而不是一般三维旋转能力不足。
\end{knowledgebox}

\begin{knowledgebox}{心理实验能证明什么}
它说明人类 reasoning 会选择并依赖 reference frame，且更换 frame 会改变可访问的关系；它不能证明神经元一定用 GLOM embedding、matrix capsule 或某个特定 attention rule。这里的证据角色是 motivation，不是 implementation identification。
\end{knowledgebox}

\subsection{Six rods：同一输入可以有完全不同的 parse}

六根杆既可被看成一个三瓣 crown，也可被看成 central rectangle 加若干斜瓣。两种 percept 对外部线段位置没有争议，却让哪些边“天然成组”、哪些边显得 parallel 完全不同。Hinton 将它类比为一句话的 structural ambiguity：truth condition 可能相同，parse 却不同。

\lecturefigure{slide-10-six-rods-percept.jpg}{六根杆的另一种 percept：同一几何输入可被不同结构组织。}{Stanford Online 官方视频 00:16:08--00:16:55。}

\lecturefigure{slide-11-alternative-percepts.jpg}{不同 internal representation 可以解释同一外部事实。}{Stanford Online 官方视频 00:16:56--00:17:55。}

\teachervoice{讲者强调这不像 Necker cube：不是你相信外部世界的 depth 改变了，而是你用另一种结构看同一事实。他用 “next weekend we should be visiting relatives” 的两种句法读法说明：representation change 不一定改变 truth condition。}

\subsection{Intrinsic relation 放进 weights，viewer pose 放进 activity}

“crown 与 flap 的固定关系”应当随 viewpoint 保持不变，适合由 learned weights 表示；“crown 相对观察者的姿态”则随视角变化，适合由 dynamic activity 表示。课堂图中的 $R_{wx}$、$R_{wv}$、$R_{xv}$ 可理解为在 object/part/viewer frames 之间的变换。

\lecturefigure{slide-12-crown-parse-tree.jpg}{Crown parse tree：边上保存 parent-child intrinsic coordinate relation。}{Stanford Online 官方视频 00:17:56--00:17:59。}

\lecturefigure{slide-13-zigzag-parse-tree.jpg}{Zig-zag/rectangle parse：相同 rods 对应另一棵结构树。}{Stanford Online 官方视频 00:18:00--00:20:41。}

\begin{knowledgebox}{读图：比较 node grouping，而不是比较线段位置}
两棵树的叶节点仍是同一组六根 rods；变化的是中间 node：crown parse 把三块 flap 组织到同一 whole，zig-zag parse 则突出 rectangle 与上下斜瓣。边上的 $R$ 表示 intrinsic frame relation。读图时先比较哪些叶子共享 parent，再比较哪些 relation 随 viewpoint 应保持不变；不要把两张图误认为两个不同外部物体。
\end{knowledgebox}

一种简化的 composition 可写为：

\[
R_{xv}=R_{wv}R_{xw}.
\]

其中，$R_{xw}$ 表示 part $x$ 相对 whole $w$ 的 intrinsic relation，$R_{wv}$ 表示 whole 相对 viewer $v$ 的当前 pose，$R_{xv}$ 是推得的 part-to-viewer pose。不同 convention 可能改变乘法顺序；教学重点是把 viewpoint-invariant knowledge 与 viewpoint-dependent state 分开。

\begin{warningbox}{Matrix equation 不是课堂实现承诺}
GLOM 论文讨论多种 pose/identity representation，并没有要求所有 relation 都必须是一张显式刚体变换矩阵。实际 embedding 还要容纳 identity、uncertainty、texture 与非刚体结构。把所有语义硬塞进 $SE(3)$ 会过度简化。
\end{warningbox}

\subsection{Mental image 不是 pixel buffer}

Hinton 的 mental imagery 观点是：内部图像可以是 structural description 加一组 viewer-relative transforms。为了回答空间问题，人会隐式选择 orientation、scale 与 position；如果换一个 frame，推理路径与“显眼关系”也会变化。

\lecturefigure{slide-14-mental-image-crown.jpg}{Mental image：在结构树上补充 node-to-viewer pose，便于传播 viewpoint。}{Stanford Online 官方视频 00:20:42--00:22:01。}

\teachervoice{课堂让听众想象“向东一英里、向北一英里、再向东一英里”。多数人会默认 north 朝上、选定某个尺度与画面位置，尽管解题并不要求这些选择。这个小实验揭示：reference frame 往往是 reasoning 的隐变量。}

\subsection{本章小结}

Part-whole hierarchy 需要同时表达 identity 与 pose。GLOM 想把稳定的 relation knowledge 放进网络函数，把输入特定的 viewpoint 与 parse 放进 embedding activity；cube、six rods 与 mental imagery 为这种分工提供认知动机，但不替代模型实验。

